\chapter{\ploren{Cytowania i bibliografia}{Citations and bibliography}}
\label{app:cytowania}

\begin{quote}\itshape
  \ploren
  {Odwołanie bibliograficzne powinno wskazywać źródło konkretnej informacji, metody, danych lub zapożyczonego elementu. Sama obecność pozycji w bibliografii nie zastępuje cytowania jej we właściwym miejscu tekstu.}
  {A bibliographic reference should identify the source of a specific statement, method, dataset or reused element. Merely listing an item in the bibliography does not replace citing it at the relevant point in the text.}
\end{quote}

\section{\ploren{Kiedy należy podać źródło}{When a source is required}}

Źródło należy wskazać przy parafrazie cudzej myśli, definicji, danych liczbowych, metodzie, modelu, wymaganiu normatywnym, fragmencie kodu oraz materiale graficznym pochodzącym od innego autora. Cytowanie jest potrzebne również wtedy, gdy treść została przetłumaczona albo przedstawiona własnymi słowami. Nie cytuje się wiedzy powszechnej, jednak w przypadku wątpliwości bezpieczniej wskazać wiarygodne źródło.

Odsyłacz powinien znajdować się możliwie blisko informacji, którą dokumentuje. Jedno cytowanie na końcu długiego akapitu może być niejednoznaczne, jeżeli akapit zawiera kilka tez pochodzących z różnych publikacji. Cytowania nie należy dodawać wyłącznie do tytułu rozdziału ani pozostawiać bez wyjaśnienia, jaki fragment źródła wykorzystano.

Pierwszeństwo powinny mieć źródła pierwotne: artykuł opisujący metodę, norma ustanawiająca wymaganie, dokumentacja producenta danego urządzenia albo oficjalna dokumentacja oprogramowania. Opracowanie przeglądowe jest przydatne do przedstawienia stanu wiedzy i odnalezienia publikacji źródłowych, lecz nie zawsze powinno być jedyną podstawą opisu.

\section{\ploren{Podstawowe polecenia cytowania}{Basic citation commands}}

Standardowe odwołanie numeryczne tworzy polecenie \verb|\cite|. Przykładowo opis metody elementów skończonych można poprzeć publikacją przeglądową~\cite{jagota2013finite}. Polecenie \verb|\textcite| włącza nazwisko autora do zdania; przykładowo \textcite{zielinski2007cyfrowe} omawia podstawy cyfrowego przetwarzania sygnałów. Kod obu wariantów przedstawia listing~\ref{lst:cytowanie-podstawowe}.

\begin{lstlisting}[style=thesis,caption={Pojedyncze cytowanie i cytowanie włączone do zdania},label={lst:cytowanie-podstawowe}]
Opis metody przedstawiono w literaturze \cite{jagota2013finite}.

\textcite{zielinski2007cyfrowe} omawia podstawy cyfrowego
przetwarzania sygnalow.
\end{lstlisting}

Przy wskazywaniu konkretnego fragmentu książki albo obszernego raportu warto podać numer strony, rozdziału lub sekcji. Opcjonalny argument polecenia jest automatycznie umieszczany przed numerem źródła, jak pokazano w listingu~\ref{lst:cytowanie-strona}.

\begin{lstlisting}[style=thesis,caption={Cytowanie ze wskazaniem strony lub rozdziału},label={lst:cytowanie-strona}]
Opis algorytmu przedstawiono w monografii
\cite[s.~125--128]{zielinski2007cyfrowe}.

Szczegoly implementacji omowiono w
\cite[rozdz.~4]{ramalho2015zaawansowany}.
\end{lstlisting}

W pracy anglojęzycznej skróty lokalizatora powinny być zapisane po angielsku, na przykład \verb|p.|, \verb|pp.| albo \verb|chap.|. Lokalizator musi odpowiadać miejscu wykorzystanemu przez autora; nie powinien być kopiowany z cytowania wtórnego bez sprawdzenia oryginału.

\section{\ploren{Kilka źródeł w jednym odwołaniu}{Multiple sources in one citation}}

Jeżeli kilka publikacji wspólnie potwierdza jedno stwierdzenie, ich klucze umieszcza się w jednym poleceniu, rozdzielając je przecinkami. Nie trzeba ręcznie porządkować kluczy według numerów bibliografii. Konfiguracja \texttt{sortcites=true} porządkuje odsyłacze, a styl \texttt{numeric-comp} łączy kolejne numery w zakres. Zapis z listingu~\ref{lst:cytowanie-wielokrotne} daje jedno zbiorcze odwołanie~\cite{knuth1984,lamport1994,mittelbach2004}; ponieważ pozycje mają kolejne numery, są wyświetlane jako zakres.

\begin{lstlisting}[style=thesis,caption={Kilka pozycji w jednym odwołaniu bibliograficznym},label={lst:cytowanie-wielokrotne}]
Zasady skladu dokumentu opisano w kilku zrodlach
\cite{knuth1984,lamport1994,mittelbach2004}.
\end{lstlisting}

Kilka źródeł należy łączyć tylko wtedy, gdy odnoszą się do tej samej tezy. Umieszczenie długiej listy publikacji po całym akapicie nie wyjaśnia, które źródło wspiera poszczególne stwierdzenia. Nie należy również tworzyć osobnych sąsiadujących odwołań \verb|\cite{...}\cite{...}|, jeżeli mogą zostać zapisane jako jedno cytowanie.

\section{\ploren{Kolejność bibliografii}{Bibliography ordering}}

Szablon stosuje styl numeryczny IEEE i ustawienie \texttt{sorting=none}. Pozycje są numerowane według kolejności pierwszego cytowania w pracy, a nie alfabetycznie. Jest to naturalny wybór dla prac technicznych: czytelnik otrzymuje małe numery w kolejności pojawiania się źródeł. Kolejne cytowanie tej samej pozycji zachowuje wcześniej nadany numer.

Nie należy ręcznie wpisywać numerów, na przykład \verb|[7]|, ani uzależniać treści zdania od bieżącej numeracji. Dodanie nowego źródła wcześniej w dokumencie może zmienić wszystkie kolejne numery, natomiast klucze w poleceniach \verb|\cite| pozostają bez zmian.

Bibliografia zawiera domyślnie tylko pozycje cytowane. Polecenie \verb|\nocite{*}| dodaje wszystkie rekordy z pliku \texttt{.bib}, lecz w pracy dyplomowej zwykle nie powinno być używane: bibliografia nie jest wykazem wszystkiego, co autor przeczytał, ale wykazem źródeł rzeczywiście przywołanych w dokumencie.

\section{\ploren{Budowa pliku bibliograficznego}{Structure of the bibliography file}}

Każdy rekord w pliku \texttt{bibliografia.bib} składa się z typu, unikatowego klucza i pól opisujących dokument. Klucz powinien być krótki, stabilny i czytelny, na przykład \texttt{nazwisko2024metoda}. Zmiana klucza wymaga poprawienia wszystkich odwołań w plikach \texttt{.tex}. Kompletny przykładowy rekord artykułu zawiera listing~\ref{lst:bib-artykul}.

\begin{lstlisting}[style=thesis,caption={Przykładowy rekord artykułu naukowego z numerem DOI},label={lst:bib-artykul}]
@article{kowalski2024diagnostics,
  author       = {Jan Kowalski and Anna Nowak},
  title        = {Diagnostics of an Electric Drive},
  journaltitle = {Journal of Electrical Engineering},
  year         = {2024},
  volume       = {75},
  number       = {2},
  pages        = {101--112},
  doi          = {10.1234/example.2024.15}
}
\end{lstlisting}

W tytule należy zachować wielkie litery istotne dla nazw własnych i skrótów, obejmując je dodatkową parą nawiasów klamrowych, na przykład \verb|{MATLAB}|, \verb|{PLC}| lub \verb|{IEEE}|. Zakres stron zapisuje się dwoma łącznikami, na przykład \verb|101--112|. W polu \texttt{doi} podaje się sam identyfikator, bez prefiksu \texttt{https://doi.org/}.

\subsection{\ploren{Książki i rozdziały}{Books and chapters}}

Dla książki podstawowe pola to \texttt{author}, \texttt{title}, \texttt{year} i \texttt{publisher}. Rozdział w pracy zbiorowej opisuje się typem \texttt{@incollection}, podając również tytuł książki, redaktorów i strony. Nie należy zastępować autora nazwą sklepu, serwisu katalogowego ani witryny, z której pobrano opis.

\subsection{\ploren{Materiały konferencyjne i prace dyplomowe}{Conference papers and theses}}

Artykuł konferencyjny zapisuje się zwykle jako \texttt{@inproceedings}, a pracę doktorską lub magisterską jako \texttt{@thesis}. Warto podać pełną nazwę konferencji lub uczelni, rok, strony, DOI oraz miejsce publikacji, jeżeli dane te są dostępne.

\subsection{\ploren{Źródła internetowe}{Online sources}}

Oficjalną dokumentację dostępną wyłącznie w sieci najlepiej zapisać jako \texttt{@online}. Pole \texttt{urldate} określa dzień dostępu w formacie ISO. Jeżeli dokument ma autora instytucjonalnego, nazwę organizacji obejmuje się dodatkową parą nawiasów klamrowych, aby BibLaTeX nie potraktował jej jak imienia i nazwiska. Przykład takiego rekordu przedstawiono w listingu~\ref{lst:bib-online}.

\begin{lstlisting}[style=thesis,caption={Przykładowy opis oficjalnej dokumentacji internetowej},label={lst:bib-online}]
@online{stmicroelectronics2026stm32,
  author  = {{STMicroelectronics}},
  title   = {{STM32} Microcontrollers},
  url     = {https://www.st.com/en/microcontrollers.html},
  urldate = {2026-07-19}
}
\end{lstlisting}

Blog, forum lub materiał promocyjny może pomóc w rozwiązaniu problemu technicznego, ale nie powinien bez uzasadnienia zastępować dokumentacji, normy albo recenzowanej publikacji. Datę dostępu podaje się szczególnie dla treści, które mogą się zmieniać.

\subsection{\ploren{Normy, oprogramowanie i dane}{Standards, software and data}}

Normy można opisywać jako \texttt{@standard}, jeżeli używana wersja BibLaTeX obsługuje ten typ, albo jako \texttt{@report}. Trzeba podać oznaczenie, pełny tytuł, organizację wydającą i rok konkretnego wydania. Samo powołanie się na „normę IEC” bez numeru i wydania jest niewystarczające.

Oprogramowanie oraz zbiory danych powinny mieć autora lub organizację, nazwę, wersję, rok, adres i trwały identyfikator, jeżeli został nadany. Preferowany jest DOI archiwalnego wydania z repozytorium, a nie odsyłacz do zmieniającej się strony głównej projektu. W informacji o dostępności kodu i danych można podać adres repozytorium, ale narzędzie lub zbiór wykorzystany jako źródło nadal powinien zostać opisany w bibliografii.

\section{\ploren{DOI, adresy URL i metadane}{DOI, URLs and metadata}}

Styl używany przez szablon wyświetla numer DOI, jeżeli rekord zawiera poprawne pole \texttt{doi}; przykładem jest artykuł dotyczący obliczeń pola magnetycznego~\cite{5424046}. DOI jest trwalszy od adresu strony wydawcy, dlatego nie należy wpisywać obu identyfikatorów, jeśli URL prowadzi wyłącznie do tej samej publikacji. Adres URL jest natomiast właściwy dla dokumentacji internetowej, repozytorium lub materiału bez DOI.

Rekord pobrany z bazy bibliograficznej trzeba sprawdzić. Typowe problemy eksportu to: nieprawidłowa wielkość liter w tytule, brak DOI, zdublowany prefiks DOI, autor zapisany w polu tytułu, niepełny zakres stron, skrócona nazwa instytucji oraz pola zawierające kod HTML. Automatyczny import oszczędza czas, ale nie zwalnia z kontroli metadanych.

\section{\ploren{Źródła rysunków, tabel i kodu}{Sources of figures, tables and code}}

Jeżeli rysunek lub tabela zostały przejęte bez zmian, cytowanie powinno znajdować się w podpisie albo w bezpośrednio poprzedzającym opisie. Przy adaptacji należy jasno zaznaczyć, że element opracowano na podstawie wskazanego źródła. Samo umieszczenie pozycji w bibliografii nie informuje, z którego materiału pochodzi element.

Kod zaczerpnięty z dokumentacji, biblioteki lub repozytorium wymaga podania źródła i sprawdzenia licencji. Krótkie, powszechne konstrukcje języka nie wymagają cytowania, lecz charakterystyczny algorytm lub większy fragment implementacji nie może być przedstawiony jako własny. Informacja o licencji może być potrzebna niezależnie od odwołania bibliograficznego.

\section{\ploren{Narzędzia do zarządzania literaturą}{Reference management tools}}

Menedżer bibliografii porządkuje opisy publikacji, pliki PDF, notatki i słowa kluczowe, a następnie eksportuje dane do formatu obsługiwanego przez LaTeX. Nie zastępuje jednak wyszukiwania literatury ani oceny jakości źródeł. Zalecany proces obejmuje zatem osobno: odnalezienie publikacji, sprawdzenie jej wiarygodności, zapisanie rekordu, kontrolę metadanych i wykorzystanie go w pliku \texttt{bibliografia.bib}.

\subsection{\ploren{Gdzie wyszukiwać literaturę}{Where to search for literature}}

Dla elektrotechniki, elektroniki, automatyki i informatyki podstawowym źródłem publikacji jest \href{https://ieeexplore.ieee.org/}{IEEE Xplore}. W informatyce istotna jest również \href{https://dl.acm.org/}{ACM Digital Library}. Zależnie od tematu przydatne są serwisy wydawców, między innymi \href{https://www.sciencedirect.com/}{ScienceDirect}, \href{https://link.springer.com/}{SpringerLink}, \href{https://onlinelibrary.wiley.com/}{Wiley Online Library} i \href{https://www.mdpi.com/}{MDPI}. MDPI jest wydawcą czasopism, a nie ogólną bazą indeksującą całą literaturę naukową.

Bazy \href{https://www.scopus.com/}{Scopus} i \href{https://www.webofscience.com/}{Web of Science} umożliwiają wyszukiwanie przekrojowe oraz śledzenie cytowań. \href{https://scholar.google.com/}{Google Scholar} jest wygodnym narzędziem do rozpoczęcia poszukiwań, odnajdywania podobnych prac i sprawdzania, kto cytował daną publikację. Wynik wyszukiwarki nie jest jednak źródłem bibliograficznym: cytować należy odnaleziony artykuł, książkę, normę albo dokumentację.

Repozytoria preprintów, na przykład \href{https://arxiv.org/}{arXiv}, pozwalają poznać najnowsze wyniki, ale zamieszczona tam wersja może nie być recenzowana. Jeżeli istnieje późniejsza wersja opublikowana w czasopiśmie lub materiałach konferencyjnych, zwykle należy odwołać się właśnie do niej. Dostęp do części baz i pełnych tekstów może wymagać zalogowania przez bibliotekę uczelni albo wykorzystania uczelnianego serwera pośredniczącego.

Wyszukiwanie warto rozpocząć od kilku zestawów słów kluczowych, przede wszystkim w języku angielskim. Po znalezieniu dobrej publikacji należy sprawdzić zarówno jej bibliografię, jak i prace, które ją cytują. Takie wyszukiwanie wsteczne i do przodu jest zwykle skuteczniejsze niż ograniczenie się do pierwszej strony wyników.

Większość baz naukowych i serwisów wydawców umożliwia pobranie opisu znalezionej publikacji. Polecenie może być oznaczone jako \emph{Cite}, \emph{Export citation}, \emph{Download citation} albo podobnie. Dla tego szablonu należy wybrać format BibTeX lub BibLaTeX, jeżeli jest dostępny, a pobrany rekord wkleić do pliku \texttt{bibliografia.bib}. Format RIS można zaimportować do Zotero lub Mendeley, a następnie wyeksportować rekord do formatu bibliograficznego używanego przez LaTeX.

Rekord najlepiej pobierać ze strony konkretnego artykułu w bazie lub serwisie wydawcy. Google Scholar również udostępnia eksport między innymi do BibTeX-u, jednak uzyskane dane mogą być skrócone albo niekompletne. Niezależnie od miejsca pobrania trzeba sprawdzić autora, tytuł, rok, nazwę czasopisma lub konferencji, tom, numer, strony i DOI. Eksport rekordu nie oznacza automatycznie, że użytkownik uzyskał dostęp do pełnego tekstu publikacji.

\subsection{\ploren{Ocena i wybór źródeł}{Evaluation and selection of sources}}

Przed wykorzystaniem publikacji należy ustalić, czy jest ona recenzowana, odpowiada tematowi i rzeczywiście potwierdza formułowane stwierdzenie. Znaczenie mają także aktualność, reputacja czasopisma lub konferencji, kompetencje autorów, kompletność metody i możliwość zweryfikowania wyników. Liczba cytowań może być pomocniczą wskazówką, lecz sama nie przesądza o jakości pracy, zwłaszcza w przypadku nowych publikacji.

Należy preferować źródła pierwotne i pełne wersje dokumentów. Warto sprawdzić DOI, autorów, afiliacje, rok, numer tomu, strony oraz informację o ewentualnym wycofaniu publikacji. Materiały promocyjne, przypadkowe strony internetowe i odpowiedzi z forów mogą pomóc w rozpoznaniu problemu, ale zazwyczaj nie powinny zastępować recenzowanych artykułów, norm, monografii ani oficjalnej dokumentacji.

\subsection{\ploren{Zotero --- rozwiązanie zalecane}{Zotero --- recommended solution}}

Zotero jest zalecanym w tym szablonie menedżerem literatury. Program pozwala tworzyć kolekcje, dołączać pliki PDF, oznaczać rekordy słowami kluczowymi, przechowywać notatki oraz wykrywać duplikaty. Rozszerzenie \emph{Zotero Connector} zapisuje publikacje i ich metadane bezpośrednio z przeglądarki. Aktualna instrukcja instalacji i obsługi znajduje się w \href{https://www.zotero.org/support/connector}{dokumentacji Zotero Connector}.

Zalecany sposób pracy z Zotero jest następujący:
\begin{enumerate}
  \item utworzyć osobną kolekcję dla przygotowywanej pracy;
  \item zapisywać publikacje z ich stron w bazach naukowych, a nie wyłącznie z widoku pliku PDF;
  \item sprawdzić autorów, tytuł, typ dokumentu, rok, nazwę czasopisma lub konferencji, strony i DOI;
  \item połączyć duplikaty i uzupełnić brakujące dane;
  \item wyeksportować wybrane rekordy do formatu BibLaTeX lub BibTeX;
  \item umieścić rekordy w pliku \texttt{bibliografia.bib} i cytować je za pomocą stabilnych kluczy.
\end{enumerate}

Automatycznie pobrane metadane zawsze wymagają kontroli. Szczególnie często trzeba poprawić wielkość liter w skrótach technicznych, kolejność imienia i nazwiska, typ publikacji, zakres stron, DOI oraz nazwę materiałów konferencyjnych. Klucze cytowań powinny być generowane według jednej reguły i nie powinny zmieniać się przy każdym eksporcie.

Przy pracy nad jednym dokumentem najlepiej utrzymywać jeden plik \texttt{bibliografia.bib}. Nie ma potrzeby tworzenia osobnej bibliografii dla każdego rozdziału. Jeżeli eksport z Zotero jest wykonywany wielokrotnie, należy uważać, aby nie nadpisać ręcznych poprawek w pliku \texttt{.bib}; bezpieczniej poprawić dane źródłowe w Zotero i ponowić eksport.

\subsection{\ploren{Zotero i Prism}{Zotero and Prism}}

Prism umożliwia połączenie projektu z indywidualną biblioteką Zotero, co ułatwia wyszukiwanie zapisanych publikacji i wstawianie cytowań podczas pracy nad dokumentem. Zakres integracji, sposób uwierzytelnienia i obsługa bibliotek grupowych mogą zmieniać się wraz z rozwojem usługi, dlatego należy korzystać z \href{https://help.openai.com/en/articles/20001050-troubleshooting-and-getting-help-in-prism}{aktualnej dokumentacji Prism}.

Integracja z edytorem nie zwalnia z kontroli pliku \texttt{bibliografia.bib}. Po przeniesieniu projektu między Prism, Overleafem i komputerem lokalnym trzeba sprawdzić, czy plik bibliograficzny został dołączony, klucze cytowań nie uległy zmianie, a wszystkie rekordy wykorzystane w pracy są dostępne również poza usługą, w której je dodano.

\subsection{\ploren{Mendeley, JabRef i inne rozwiązania}{Mendeley, JabRef and other solutions}}

Mendeley Reference Manager jest alternatywnym menedżerem literatury, a rozszerzenie \emph{Mendeley Web Importer} pozwala zapisywać rekordy z przeglądarki. Aktualne instrukcje są dostępne w \href{https://www.mendeley.com/guides}{przewodnikach Mendeley}. JabRef jest natomiast narzędziem skoncentrowanym bezpośrednio na plikach BibTeX i BibLaTeX, dlatego dobrze pasuje do lokalnego sposobu pracy z LaTeX-em.

Nie należy równolegle utrzymywać kilku niezależnych bibliotek tej samej pracy bez ustalenia, która z nich jest źródłem nadrzędnym. Niezależnie od wybranego programu końcowym interfejsem szablonu pozostaje poprawny i przenośny plik \texttt{bibliografia.bib}.

\subsection{\ploren{Zalecany przepływ pracy}{Recommended workflow}}

Cały proces można podsumować następująco:
\begin{center}
  \small
  wyszukiwanie $\rightarrow$ ocena źródła $\rightarrow$ zapis w menedżerze
  $\rightarrow$ kontrola metadanych $\rightarrow$ eksport pliku \texttt{.bib}
  $\rightarrow$ cytowanie w LaTeX-u.
\end{center}

Przed oddaniem pracy należy ponownie wyszukać brakujące pola, zduplikowane rekordy oraz pozycje różniące się tylko sposobem zapisu DOI lub nazwiska autora. Warto również zachować kopię biblioteki menedżera i pliku \texttt{bibliografia.bib} razem z pozostałymi materiałami źródłowymi pracy.

\section{\ploren{Najczęstsze błędy}{Common mistakes}}

Do najczęstszych błędów należą: ręczne wpisywanie numerów źródeł, cytowanie tylko na końcu całego rozdziału, brak stron przy dokładnym przytoczeniu, powoływanie się na źródło wtórne mimo dostępności oryginału, umieszczanie niecytowanych pozycji przez \verb|\nocite{*}|, niekompletne dane internetowe, cytowanie wyniku wyszukiwarki zamiast dokumentu, niespójne klucze oraz brak kontroli rekordów zaimportowanych automatycznie.

Przed złożeniem pracy należy również sprawdzić, czy każdy odsyłacz prowadzi do właściwej pozycji, wszystkie cytowane rekordy są widoczne w bibliografii, DOI i URL są aktywne, nazwiska i tytuły nie zawierają błędów oraz czy kolejność numerów odpowiada pierwszemu wystąpieniu źródeł w tekście.
